\section{Conclusions}
\label{sec:conclusion}

Architecture exploration for heterogeneous distributed cyber-physical systems is limited
less by the size of the design space than by the cost of evaluating a single design point.
This paper presented HCA-DSE, a hierarchical constraint-aware exploration method that
removes structurally ill-formed, provably infeasible and provably dominated regions of the
space before any full evaluation is performed. Its technical core is a family of
admissible bounds for \emph{contention-dependent} objectives, obtained by a term-wise
relaxation combined with a decision ordering under which the queueing term of a tier
becomes exactly computable on a partial design. This removes the assignment-monotonicity
assumption on which existing exact exploration methods rely, and we proved that the exact
Pareto front is preserved.

Across four design-space scales, three workload classes and spaces of up to
$8.3\times10^{7}$ candidate architectures, HCA-DSE performed 26 to 61 full evaluations,
a reduction of $99.90\%$ to $99.98\%$ with respect to the raw space, while reproducing the
exhaustive Pareto front exactly wherever exhaustive enumeration is tractable. Both the
admissibility of the bounds and the equality of the fronts were verified exhaustively by
machine. Budget-matched NSGA-II required $1000$ evaluations to reach a median Pareto
recall of $0.50$--$0.91$. The break-even evaluator cost is in the microsecond range, and
with a discrete-event evaluator the measured speed-up reaches $558\times$. On a
containerised edge--gateway--cloud testbed with per-node link shaping, 36 architectures
and 108 measured runs, the analytical ordering agreed with the measured ordering in
$98\%$ of the architecture pairs that the model separates by more than $25\%$.

The method is intended as design-time decision support: it does not require runtime
adaptation, learned performance models, or a high-fidelity simulator in the inner loop.
Its natural continuation is the combination with an evaluator-centric methodology, in
which analytical bounds reduce an industrial design space to a few dozen candidates and a
trace-derived simulator evaluates only those; and the derivation of tighter, possibly
surrogate-assisted bounds that remain provably admissible.
